% from ega4-18-fr.tex
\oldpage[IV-4]{203}

qui vérifient cette condition engendrant l'algèbre symétrique $\mathbf S_A(E)$); changeant au besoin de base dans $E$, on peut donc supposer que $t=T_0$, et $u(T_i)=0$ pour $1\leq i\leq m-1$, de sorte que $S_{(t)}$ s'identifie canoniquement à $A[T_1,\ldots,T_{m-1}]$. Les éléments de $H$ sont alors les éléments de la forme
\[
  \sum_\alpha L_\alpha(\mathbf T)\otimes\mathbf T^\alpha
\]
de $E\otimes_A A[T_1,\ldots,T_{m-1}]$ où les $L_\alpha(\mathbf T)=\lambda_{\alpha0}+\lambda_{\alpha1}T_1+\cdots+\lambda_{\alpha,m-1}T_{m-1}$ doivent être tels que $\sum_\alpha L_\alpha(\mathbf T)\mathbf T^\alpha=0$ dans $A[T_1,\ldots,T_{m-1}]$. On vérifie aisément que lorsque la condition $\lambda_{\mathbf0,0}=0$ est remplie, les valeurs des $\lambda_{\alpha0}$ peuvent être prises \emph{arbitrairement} pour $|\alpha|\geq1$, et on peut toujours alors déterminer les $\lambda_{\alpha i}$ pour $i\geq1$ de façon à vérifier la condition $\sum_\alpha L_\alpha(\mathbf T)\mathbf T^\alpha=0$. Les éléments de l'idéal $\mathscr K$ correspondants sont alors les polynômes $\sum_\alpha\lambda_{\alpha0}\mathbf T^\alpha$ \emph{sans terme constant}, et comme $(\mathbf S_A(A))_{(t)}$ s'identifie à $A$, ces polynômes sont exactement les éléments du noyau de l'homomorphisme $(\mathbf S_A(E))_{(t)}\to(\mathbf S_A(A))_{(t)}$. On a ainsi prouvé (19.2.1) que l'immersion $j:X'\to P$ est \emph{transversalement régulière} relativement à $X$ en tous les points de $X'-Y'$ et que $j(X')=Z$ en ces points.

Reste à examiner ce qui se passe au-dessus d'un point $x\in Y$. L'hypothèse de « régularité » de $u$ en ce point faite dans (iii) équivaut (en vertu de (19.1.1) et de ($0$, 17.1.7)) à dire que $X$ est \emph{régulier} au point $x$ et que l'homomorphisme
\[
  u_x\otimes1:\mathscr E_x\otimes_{\sh{O}_{X,x}}\mathbf k(x)\longrightarrow\mathfrak m_x/\mathfrak m_x^2
\]
est \emph{injectif}. Cela étant, comme le morphisme $p:P\to X$ est lisse (17.3.9), pour tout $z\in P$ au-dessus de $x$, $P$ est régulier au point $z$ (17.5.8); en outre ($0$, 17.3.3) l'homomorphisme canonique
\begin{equation}
  (\mathfrak m_x/\mathfrak m_x^2)\otimes_{\mathbf k(x)}\mathbf k(z)\longrightarrow\mathfrak m_z/\mathfrak m_z^2
  \tag{19.4.11.1}\label{IV.19.4.11.1}
\end{equation}
est \emph{injectif}. On en déduit que l'homomorphisme
\[
  (\mathscr E_x\otimes_{\sh{O}_{X,x}}\sh{O}_{P,z})\otimes_{\sh{O}_{P,z}}\mathbf k(z)\longrightarrow\mathfrak m_z/\mathfrak m_z^2
\]
déduit de $p^*(u)$ est aussi injectif, car il s'écrit comme le composé
\[
  (\mathscr E_x\otimes_{\sh{O}_{X,x}}\mathbf k(x))\otimes_{\mathbf k(x)}\mathbf k(z)\longrightarrow(\mathfrak m_x/\mathfrak m_x^2)\otimes_{\mathbf k(x)}\mathbf k(z)\longrightarrow\mathfrak m_z/\mathfrak m_z^2
\]
où la première flèche est injective par platitude et la seconde est l'homomorphisme injectif (19.4.11.1). Comme $\mathscr H$ est (dans un voisinage de $z$ dans $P$) facteur direct de $\mathscr E\otimes_{\sh{O}_X}\sh{O}_P=p^*(\mathscr E)$, l'homomorphisme $v_z\otimes1:\mathscr H_z\otimes_{\sh{O}_{P,z}}\mathbf k(z)\to\mathfrak m_z/\mathfrak m_z^2$ est aussi injectif, ce qui établit la « régularité » de l'homomorphisme $v$ au point $z$. Reste à voir que l'image par $j$ de $X'$ est identique au sous-préschéma fermé $Z$. Il suffira pour cela de montrer d'une part que l'image par $j$ des points de $q^{-1}(x)$ (où $q:X'\to X$ est le morphisme structural) sont exactement ceux de $Z\cap p^{-1}(x)$ et d'autre part qu'en chacun de ces points $z=j(x')$, l'homomorphisme surjectif $\sh{O}_{Z,z}\to\sh{O}_{X',x'}$ est bijectif. Mais, compte tenu de l'hypothèse et de (19.4.3), la restriction $Y'\to Y$ de $q$ est lisse, et puisque $\sh{O}_{Y,x}$ est supposé régulier, il en est de même de $\sh{O}_{Y',x'}$ (17.5.8); en outre, l'immersion $Y'\to X'$ est régulière (19.4.2), donc $\sh{O}_{X',x'}$ est aussi un anneau régulier (19.1.1). Pour prouver que l'homomorphisme $\sh{O}_{Z,z}\to\sh{O}_{X',x'}$ est bijectif, il suffit donc de montrer que les

